\documentclass[11pt,a4paper]{article}
\usepackage[margin=25mm]{geometry}
\usepackage{amsmath,amssymb,amsthm}
\newtheorem{theorem}{Theorem}
\newtheorem{proposition}[theorem]{Proposition}
\newtheorem{corollary}[theorem]{Corollary}
\newcommand{\R}{\mathbb R}
\newcommand{\one}{\mathbf 1}
\newcommand{\im}{\operatorname{im}}
\newcommand{\Stab}{\operatorname{Stab}}
\newcommand{\Aut}{\operatorname{Aut}}
\title{Realization cones from a finite vertex action}
\author{Human Compression project}\date{}
\begin{document}\maketitle

The realization theory discussed in Section 6 of Schulte--Weiss~\cite{survey}
has a common linear algebraic description. Once the action on vertices is
fixed, its realizations are described by positive semidefinite matrices
commuting with that action. We give the argument, including the distinction
between an equivariant placement, a faithful realization of faces, and a
realization with a prescribed full symmetry group. The cone and its
representation-theoretic interpretation are established parts of realization
theory, attributed in the survey to McMullen and McMullen--Monson. They are
not new theorems claimed by this project.

\section{Distances and centered Gram matrices}
Let $X$ be a finite nonempty set of cardinality $m$, and let a finite group $G$ act on $X$.
Write $M=\R^X$ with its usual orthonormal basis $(e_x)_{x\in X}$, and let
$P_g e_x=e_{gx}$. Put
\[
 J=I-\frac1m\one\one^{\mathsf T},\qquad
 M_0=\one^\perp.
\]
A placement is a family $f=(f_x)_{x\in X}$ of points in a Euclidean space.
We always replace the ambient space by the affine span of these points,
and identify placements differing by an isometry that preserves the labels.
Subtracting the centroid makes $\sum_x f_x=0$. The affine span is then the
linear span, denoted $W$.

\begin{theorem}\label{gram}
The centered placements for which every $g\in G$ is induced by an isometry
are classified, up to labeled congruence, by the cone
\[
 \mathcal C_G=\{B\in\R^{X\times X}: B=B^{\mathsf T},\ B\succeq0,
       \ B\one=0,\ P_g^{\mathsf T}BP_g=B\text{ for every }g\in G\}.
\]
The dimension of the placement is $\operatorname{rank}B$. Its squared
distance matrix is
\begin{equation}\label{distance}
 D_{xy}=B_{xx}+B_{yy}-2B_{xy},
 \qquad B=-\tfrac12 JDJ.
\end{equation}
These formulas give mutually inverse linear maps between $\mathcal C_G$
and the invariant squared distance matrices.
\end{theorem}
\begin{proof}
For a placement define $T:M\to W$ by $Te_x=f_x$. Its Gram matrix is
$B=T^*T$. It is positive semidefinite and symmetric, and $T\one=0$ gives
$B\one=0$. An isometry permuting the points preserves their centroid, so
after centering it is orthogonal. Its preservation of inner products is
exactly the displayed invariance of $B$.

Conversely let $B\in\mathcal C_G$. The positive semidefinite square root
$B^{1/2}$ commutes with every $P_g$, since each eigenspace of $B$ is
$G$-invariant. Set $W=\im B$ and $f_x=B^{1/2}e_x$. These vectors span $W$,
their sum is zero, and $g$ acts on them by the orthogonal restriction of
$P_g$. Their Gram matrix is $B$.

If two spanning families $f_x$ and $f'_x$ have this Gram matrix, the rule
$\sum_x a_xf_x\mapsto\sum_x a_xf'_x$ is well-defined: the squared norm of
either sum is $a^{\mathsf T}Ba$. It is an onto linear isometry. Thus the
matrix determines the labeled placement, including degenerate placements.
Its rank is the dimension of the span.

The first formula in~\eqref{distance} follows by expanding
$\|f_x-f_y\|^2$. Multiplying on both sides by $J$ annihilates the two
terms constant in one coordinate, and $JBJ=B$, which proves the second
formula. To see that no additional condition on a proposed distance matrix
has been hidden, take any symmetric hollow matrix $D$ and set
$B=-JDJ/2$. If $r_x=m^{-1}\sum_yD_{xy}$ and
$s=m^{-2}\sum_{x,y}D_{xy}$, then
\[
 B_{xy}=-\tfrac12(D_{xy}-r_x-r_y+s),\qquad B_{xx}=r_x-s/2.
\]
Consequently $B_{xx}+B_{yy}-2B_{xy}=D_{xy}$. If $B\succeq0$, the
construction already given realizes $D$. Invariance passes in both
directions because $J$ commutes with permutation matrices.
\end{proof}

In particular, one may specify $D$ by one coordinate for each $G$-orbit
on unordered pairs of distinct vertices. These are the diagonal classes
used in the survey. The conditions in Theorem~\ref{gram} show directly
that the set of possible diagonal vectors is a closed convex cone: it is
linearly isomorphic to the intersection of the positive semidefinite cone
with a linear subspace. No assertion that all of its points are faithful
is involved.

If $f$ and $h$ are two centered placements, their blend is the placement
$x\mapsto(f_x,h_x)$ in $W_f\oplus W_h$, restricted to the span of these
vectors. Its Gram matrix is $B_f+B_h$, and its squared distance matrix is
$D_f+D_h$. The restriction to the actual span matters when the two
representations contain equivalent summands. Multiplication of the
matrices by $c\geq0$ corresponds to multiplication of the placement by
$\sqrt c$.

\section{Irreducible realizations and fixed vectors}
\begin{theorem}\label{extreme}
A nonzero $B\in\mathcal C_G$ spans an extreme ray of $\mathcal C_G$ if and
only if the real $G$-module $\im B$ is irreducible. Every element of the
cone is a sum of matrices on extreme rays.
\end{theorem}
\begin{proof}
First suppose $W=\im B$ is irreducible and $B=C+D$ with
$C,D\in\mathcal C_G$. If $v\in\ker B$, positivity gives
$0=\langle Cv,v\rangle+\langle Dv,v\rangle$, so both terms vanish.
Diagonalizing a positive semidefinite matrix shows that a zero value of
its quadratic form puts $v$ in its kernel. Thus $C$ and $D$ vanish on
$\ker B$, and symmetry makes them preserve $W=(\ker B)^\perp$.
Every eigenspace of the symmetric operator $C|_W$ is $G$-invariant.
Irreducibility therefore makes $C|_W$ a scalar operator. The same holds
for $B|_W$, whose scalar is positive. It follows that $C$ is a nonnegative
multiple of $B$; the same is true of $D$. Hence the ray is extreme.

Conversely suppose $W$ has a nonzero proper invariant subspace $U$.
Its orthogonal complement is invariant, since the action is orthogonal.
Let $Q$ be orthogonal projection onto $U$, extended by zero on $W^\perp$.
Then $Q$ commutes with $G$. With $A=B^{1/2}$, the decomposition
\[
 B=AQA+A(I-Q)A
\]
consists of two members of $\mathcal C_G$. Both are nonzero, since $A$
is invertible on $W$, and their ranks are respectively $\dim U$ and
$\dim W-\dim U$. Neither is proportional to $B$. Thus its ray is not
extreme. Notice that $B$ need not commute with an arbitrarily chosen
invariant projection; the square roots on both sides of $Q$ avoid that
unjustified assumption.

For the final assertion use the placement $f_x$ constructed in
Theorem~\ref{gram}. Its orthogonal representation on $W$ splits into a
finite orthogonal sum of irreducible invariant subspaces: choose a
minimal nonzero invariant subspace, take its orthogonal complement,
and continue by induction on dimension. Projecting every $f_x$ onto
each summand gives centered equivariant placements. Their Gram matrices
sum to $B$, and each nonzero one has that irreducible summand as its
span, because the original $f_x$ span $W$. The first assertion applies.
\end{proof}

Now assume that the action on $X$ is transitive. Choose $x_0\in X$ and
put $H=\Stab_G(x_0)$, so that $X$ is identified with the left cosets
$G/H$. For an orthogonal representation $\rho:G\to O(V)$ write
$V^H=\{v:\rho(h)v=v\text{ for every }h\in H\}$.

\begin{proposition}\label{wythoff}
Equivariant placements of $G/H$ in $V$ are precisely
\[
 f_{gH}=\rho(g)v\qquad(v\in V^H).
\]
The placement is injective exactly when $\Stab_G(v)=H$. Its centroid is
the projection of $v$ onto $V^G$. If $V$ is nontrivial irreducible and
$0\ne v\in V^H$, the placement is centered, spans $V$, and gives an
extreme ray of $\mathcal C_G$.
\end{proposition}
\begin{proof}
The displayed expression is independent of the representative of $gH$
exactly when $v$ is fixed by $H$. Equivariance forces every placement to
have this form, by taking $v=f_{H}$. Two cosets have the same image
precisely when their quotient lies in $\Stab_G(v)$, which proves the
injectivity assertion.

The mean of the orbit is $|G|^{-1}\sum_{g\in G}\rho(g)v$: every coset
appears $|H|$ times in this sum. The averaging operator is self-adjoint
and idempotent, and has image $V^G$, hence is the orthogonal projection
onto that subspace. For a nontrivial irreducible representation $V^G=0$,
and the invariant span of the orbit of a nonzero vector is $V$. The
last assertion follows from Theorem~\ref{extreme}.
\end{proof}

This identifies the Wythoff space with a fixed-vector space. Different
irreducible representations need not contribute, since $V^H$ may be zero.
Within a representation, different fixed vectors need not give the same
ray. Thus an enumeration of irreducible characters alone is not a
description of the realization cone. No assumption about the real,
complex or quaternionic type of an irreducible representation was used.
In particular, this argument does not decide Problem 23 of the survey.

\section{Faces and full geometric symmetry}
Let $\mathcal P$ be a finite abstract polytope with vertex set $X$.
A vertex placement induces the recursive face images used in the survey:
the image of a face of rank $j\geq1$ is the set of images of its incident
faces of rank $j-1$. Consider the following condition on the abstract
polytope:
\begin{equation}\label{faces}
 F\ne F'\text{ of rank }j\quad\Longrightarrow\quad
 \{E:\operatorname{rank}E=j-1,\ E<F\}
 \ne\{E:\operatorname{rank}E=j-1,\ E<F'\}.
\end{equation}
It is required for every $j\geq1$ for which there are faces to compare.

\begin{proposition}\label{faithful}
The recursive realization is faithful on every rank if and only if the
vertex placement is injective and~\eqref{faces} holds. Subject to this
combinatorial condition, faithfulness for $B\in\mathcal C_G$ is exactly
$D_{xy}>0$ for distinct vertices $x,y$.
\end{proposition}
\begin{proof}
If two distinct faces have the same set of incident faces of the preceding
rank, their images agree for every vertex placement. Faithfulness therefore
requires~\eqref{faces}, as well as injectivity on vertices. Conversely,
suppose the images in rank $j-1$ are distinct. Equality of two sets of
such images then implies equality of the two underlying sets of faces.
Condition~\eqref{faces} makes the rank-$j$ images distinct. Induction from
rank zero proves the assertion. Finally $D_{xy}=\|f_x-f_y\|^2$, so its
strict positivity is exactly vertex injectivity.
\end{proof}

Suppose~\eqref{faces} holds and the placement is faithful. Put
$A=\Aut(\mathcal P)$, acting on $X$. This action is faithful: an
automorphism fixing every vertex fixes each succeeding rank by
\eqref{faces}. The realized faces retain their ranks and incidence.

\begin{proposition}\label{symmetry}
The full symmetry group of the faithful realization, in its affine span,
is naturally isomorphic to
\[
 A_B=\{a\in A:P_a^{\mathsf T}BP_a=B\}.
\]
Consequently, for $G\leq A$, a faithful equivariant placement has full
geometric group exactly $G$ if and only if $B\in\mathcal C_G$ and no
$a\in A\setminus G$ preserves $B$.
\end{proposition}
\begin{proof}
A symmetry of the realization permutes its faces preserving ranks and
incidence. Faithfulness gives an automorphism of $\mathcal P$, and the
symmetry preserves all squared distances. By~\eqref{distance} it
preserves $B$. Conversely, if $a\in A$ preserves $B$, the same argument
with linear relations used in Theorem~\ref{gram} shows that
$f_x\mapsto f_{ax}$ extends to an orthogonal isometry of the centered
span. The recursive definition of face images makes it act on every
face as $a$. Both constructions are unique because the vertices span.
They preserve composition and are inverse to one another.
\end{proof}

The distinction in Proposition~\ref{symmetry} is necessary for chiral
realizations. A group representation and a choice of initial vertex may
have more symmetry than the designated group. The proposition supplies
an exact criterion for that issue for a fixed finite polytope; it asserts
no existence or classification result for all polytopes.

\section{Use in the source disposition}
Theorem~\ref{gram} proves the finite cone and diagonal-vector statements
in the survey, and identifies blending with addition. Theorem~\ref{extreme}
proves the statement about pure realizations and irreducible orthogonal
representations. Proposition~\ref{wythoff} supplies the fixed-space
interpretation, with the stabilizer condition that distinguishes an
injective orbit from a merely faithful representation.
Propositions~\ref{faithful} and~\ref{symmetry} state the further obligations
needed for faithful faces and the full geometric group.

These arguments do not compute the cones for every regular toroid, classify
discrete realizations in dimension three, decide the quaternionic question,
or exhibit a faithful chiral realization of a regular abstract polytope.
They do not replace the separately attributed geometric classifications
in the survey. The present manuscript rederives common established theory;
there is no new conjecture, Lean or compression-measurement credit.

\begin{thebibliography}{9}\small
\bibitem{survey} E. Schulte and A. Ivi\'c Weiss,
\emph{Problems on polytopes, their groups, and realizations},
Period. Math. Hungar. 53 (2006), 231--255; fixed arXiv math/0608397v1,
Section 6 and its references to P. McMullen,
\emph{Realizations of regular polytopes}, Aequationes Math. 37 (1989),
38--56, and P. McMullen and B. R. Monson,
\emph{Realizations of regular polytopes, II}, Aequationes Math. 65 (2003),
102--112. Attribution here follows the read survey; its cited classifications
are not re-established by the matrix arguments above.
\end{thebibliography}
\end{document}
